% =====================================================================
% 11. Conclusion
% =====================================================================
\section{Conclusion}
\label{sec:conclusion}

Holding default probabilities fixed, the $99.9\%$ VaR of the portfolio ranges from USD~\val{q8.M2.VaR_ref|mn2} million under Gaussian factors to USD~\val{q8.M1.VaR_ref|mn2} million under the empirical distribution, with the $t$-copula model only \val{q8.M3.VaR_vs_M2|pct1}\% above the Gaussian one. The answer to our research question is therefore that the tail of the loss is sensitive to the factor distribution, but to its margins far more than to its copula: allowing the two indices to crash together matters little as long as each of them is Gaussian and the firm-specific shocks are Gaussian and independent. A natural extension, which the assignment does not ask for, would combine a $t$-copula with fat-tailed margins, and so obtain the tail thickness of M1 without tying the $99.9\%$ quantile to two historical weeks.

The same message recurs in the other parts of the analysis. Through time, the parametric models are stable because the thresholds are recalibrated, while the empirical model inherits the ghost effects of historical simulation. Against the regulatory benchmark, the model-based unexpected loss is \val{q10.ratio_UL_M2|f1} times the IRB capital, and the gap is a matter of correlation levels, not of granularity or of the shape of the factor distribution. And the storm overlay, which acts on recoveries and is independent of the credit cycle by assumption, raises the expected loss by \val{q11.dEL_exact|ppm} parts per million and ES by about USD~\val{q11.first_order.M2.addon_tail_usd|usd2}, four orders of magnitude below the Monte Carlo error of the estimates. To first order only the expected damage at the locations of the defaulting names matters, so the geographical concentration of the portfolio is irrelevant under independence; it becomes relevant only when storm years and credit years are dependent, and even then the effect stays below $0.02\%$ of ES for Swiss winter storms.

Three caveats apply. The empirical model depends on a handful of observations. With $10{,}000$ simulations, the Monte Carlo error at the $99.9\%$ level is of the same order as some of the differences of interest, which is why our conclusions rest on a reference run of two million scenarios. And the comparison with the IRB formula is driven by the systematic weights $\lambda_k$ supplied with the synthetic portfolio, which are high relative to the regulatory correlations; with weights in the regulatory range the model-based and the regulatory capital would be much closer.
